
\begin{figure}[htbp]
\centering
\begin{tikzpicture}[font=\small, x=0.8cm, y=0.65cm]
  \draw[->] (0,0) -- (9.5,0) node[right] {$p$};
  \draw[->] (0,0) -- (0,6.3) node[above] {$S_p$};
  \foreach \y/\label in {1/1,2/2,3/3,4/4,5/5,6/6} {
    \draw[black!20] (0,\y) -- (9,\y);
    \node[anchor=east] at (-0.1,\y) {\label};
  }
  \foreach \x/\label in {1/1,2/2,4/4,6/8,8/16} {
    \draw[black!20] (\x,0) -- (\x,6);
    \node[anchor=north] at (\x,-0.08) {\label};
  }
  \draw[thick, accent] plot[smooth] coordinates {(1,1) (2,1.6) (3,2.0) (4,2.28) (5,2.5) (6,2.66) (7,2.8) (8,2.91) (9,3.0)};
  \node[accent, anchor=west] at (9.1,3.0) {$f=0{,}75$};
  \draw[thick, orange!80!black] plot[smooth] coordinates {(1,1) (2,1.82) (3,2.5) (4,3.08) (5,3.57) (6,4.0) (7,4.37) (8,4.71) (9,5.0)};
  \node[orange!80!black, anchor=west] at (9.1,5.0) {$f=0{,}90$};
  \draw[dashed] (0,4) -- (9,4);
  \node[align=center] at (4.8,-1.0) {Minél nagyobb a nem párhuzamosítható rész, annál hamarabb telítődik a gyorsítás.};
\end{tikzpicture}
\caption{Amdahl törvényének szemlélete: a gyorsítás felső korlátos}
\label{fig:zv07_amdahl}
\end{figure}
